% =============================================================================
% Abstract and Keywords
% =============================================================================
\chapter*{Abstract}
\addcontentsline{toc}{chapter}{Abstract}

This report presents \eduspace{}, a web-based classroom management platform designed and developed as an end-of-year project for the 4\textsuperscript{th}-year Software Engineering program.
Built on a microservices architecture, \eduspace{} addresses the organizational and communication shortcomings of existing platforms such as Google Classroom by providing structured content management through typed posts and chapters, integrated real-time group chat, full-text search, an assignment and grading workflow, and a real-time notification system.

The system comprises seven independent backend services (User, Class, Content, File, Communication, Notification, and Search), an API Gateway, and a React single-page application.
Each service owns its dedicated PostgreSQL database and communicates asynchronously via RabbitMQ or synchronously through REST calls.
The infrastructure is containerized using Docker Compose and includes Redis for caching and rate limiting, MinIO for S3-compatible file storage, Elasticsearch for full-text search, and NGINX as a reverse proxy.

The report covers the full software development lifecycle: requirements analysis, UML modeling, system design, and implementation details including the technology stack, security measures, and DevOps practices.

\medskip
\noindent\textbf{Keywords:} microservices, classroom management, React, Node.js, Docker, WebSocket, REST API
